% META: {"id":"1SPE-TRIGO-FR-R2","chapitre":"1SPE-TRIGONOMETRIE","type_objet":"remediation","status":"approved","prerequis_testes":["R2"],"origin":"generated","status_history":["generated","audited","accepted","approved"],"release_acceptance":"RELEASE_OWNER_BATCH_ACCEPTANCE","acceptance_closure_digest":"sha256:9b3ccf9a81c5520fbb7e03b7b2d3e7bf2057a02834b6d908bf3be5f49f3b553f","acceptance_receipt_digest":"sha256:4fad86f0282e0fa4e0419cca1ad6379ae7af5a280c04a4a90880f90a1c044242"}

%---------------------------------------------------------------
% FICHE DE REMISE A NIVEAU — R2 : Reperage dans le plan
%---------------------------------------------------------------

\subsection*{Rappel — Coordonnées et distance}

Dans un repere orthonorme $(O;\vec{\imath},\vec{\jmath})$, un point $M$ a pour coordonnées $(x;y)$.

La distance $OM$ est $OM = \sqrt{x^2 + y^2}$.

\begin{exercice}{1SPE-TRIGO-FR-R2-EX1}{1}{5}
\begin{enumerate}
  \item Placer les points $A(1;0)$, $B(0;1)$, $C(-1;0)$, $D(0;-1)$ dans un repere.
  \item Calculer $OA$, $OB$, $OC$, $OD$.
  \item Ces quatre points sont-ils sur un même cercle ? Lequel ?
\end{enumerate}
\end{exercice}

% BEGIN-VERIFY
% from sympy import *
% assert sqrt(1**2 + 0**2) == 1
% assert sqrt(0**2 + 1**2) == 1
% END-VERIFY

\begin{corrige}{1SPE-TRIGO-FR-R2-EX1}
\textbf{1.} Points places sur les axes.

\textbf{2.} $OA = OB = OC = OD = 1$.

\textbf{3.} Oui, ils sont sur le cercle de centre $O$ et de rayon $1$ : le cercle trigonométrique.
\end{corrige}
